% 海洋行星（综述）
% license CCBYSA3
% type Wiki

本文根据 CC-BY-SA 协议转载翻译自维基百科\href{https://en.wikipedia.org/wiki/Ocean_world}{相关文章}。

\begin{figure}[ht]
\centering
\includegraphics[width=6cm]{./figures/76246827a4922672.png}
\caption{地球表面主要由海洋覆盖，海洋占据了地球表面的 $75\%$。因此，地球可以被视为一个水世界，但并非一个完全的海洋行星。} \label{fig_Ocwo_1}
\end{figure}
海洋世界（ocean world）、海洋行星（ocean planet）或水世界（water world）是一类行星或天然卫星，其水圈中包含大量以海洋形式存在的水，这些海洋可以位于表面之下，形成地下海洋，或直接存在于表面，甚至可能淹没所有陆地$^{[1][2][3][4]}$。“海洋世界”一词有时也被用于指代其海洋由其他流体或造海物质（thalassogen）构成的天体$^{[5]}$，例如熔岩（如木卫一 Io）、氨（与水形成共晶混合物，可能存在于土卫六 Titan 的内部海洋中），或碳氢化合物（如土卫六表面的情况，其可能是系外海洋中最常见的一类）$^{[6]}$。对地外海洋的研究被称为行星海洋学（planetary oceanography）。

地球是目前唯一已知在其表面存在液态水体的天文天体，尽管在木星的卫星欧罗巴（Europa）与盖尼米德（Ganymede），以及土星的卫星恩克拉多斯（Enceladus）与土卫六（Titan）上，人们推测其表面之下可能存在地下海洋$^{[7]}$。已经发现若干系外行星具备支持液态水存在的适宜条件$^{[8]}$。此外，在地球上也发现了大量地下水，主要以含水层的形式存在$^{[9]}$。对于系外行星而言，受限于当前技术，尚无法直接观测其表面的液态水，因此通常以大气中的水汽作为替代性指标$^{[10]}$。海洋世界的性质为理解其自身历史以及整个太阳系的形成与演化提供了重要线索；与此同时，它们在生命起源与维持方面的潜力也引起了广泛关注。

2020 年 6 月，NASA 科学家基于数学建模研究报告称，在银河系中拥有海洋的系外行星很可能是普遍存在的$^{[11][12]}$。
\subsection{概述}  
\subsubsection{定义}  
根据 Lunine 的定义，“海洋”被界定为“稳定存在、环绕整个天体的液态水体”$^{[13]}$。此外，“海洋世界”这一术语用于指代太阳系中那些拥有稳定、覆盖全球的液态水体的天体；与之相对，“海洋行星”和“水世界”则专指系外行星（即绕其他恒星运行的行星），其整体组成中含有较高质量分数的水$^{[13]}$。
\subsubsection{太阳系行星天体}
\begin{figure}[ht]
\centering
\includegraphics[width=8cm]{./figures/4b85d59dce5a977e.png}
\caption{土卫二（Enceladus）内部结构示意图} \label{fig_Ocwo_2}
\end{figure}
海洋世界因其可能孕育生命并维持生物活动的潜力而受到天体生物学家的高度关注$^{[4][3]}$。太阳系中被认为拥有地下海洋的主要卫星和矮行星尤为重要，因为与远在数光年之外、受当前技术所限难以直接探测的系外行星不同，这些天体可以通过空间探测器抵达并开展实地研究。除地球之外，太阳系中证据最为充分的水世界包括卡利斯托（Callisto）、恩克拉多斯（Enceladus）、欧罗巴（Europa）、盖尼米德（Ganymede）以及土卫六（Titan）$^{[3][14]}$。其中，欧罗巴和恩克拉多斯由于其外壳较薄且存在冰火山活动，被视为极具吸引力的探测目标。

太阳系中的其他天体也被视为可能拥有地下海洋的候选者，这一判断基于单一类型的观测证据或理论建模结果，包括天卫一（Ariel）$^{[14]}$、天卫三（Titania）$^{[15][16]}$、天卫二（Umbriel）$^{[17]}$、谷神星（Ceres）$^{[3]}$、土卫四（Dione）$^{[18]}$、土卫一（Mimas）$^{[19][20]}$、天卫五（Miranda）$^{[14]}$、天卫四（Oberon）$^{[4][21]}$、冥王星（Pluto）$^{[22]}$、海卫一（Triton）$^{[23]}$、厄里斯（Eris）$^{[4][24]}$以及鸟神星（Makemake）$^{[24]}$。
\subsubsection{系外行星}
\begin{figure}[ht]
\centering
\includegraphics[width=8cm]{./figures/d53b32c9e8bb7ae8.png}
\caption{一组含水、大小各异的系外行星，与地球进行对比（艺术家概念图；2018 年 8 月 17 日）$^{[25]}$} \label{fig_Ocwo_3}
\end{figure}
\begin{figure}[ht]
\centering
\includegraphics[width=8cm]{./figures/c7e845b1bfbd19f5.png}
\caption{以纯海洋世界作为过渡群体的系外行星族群示意图：它们位于气态巨行星与熔岩行星或岩石行星之间。} \label{fig_Ocwo_4}
\end{figure}
在太阳系之外，被描述为**候选海洋世界**的系外行星包括 GJ 1214 b $^{[26][27]}$、Kepler-22b、Kepler-62e、Kepler-62f $^{[28][29][30][31]}$，以及 Kepler-11 $^{[32]}$ 和 TRAPPIST-1 $^{[33][34]}$ 行星系统中的行星。

近年来，系外行星 TOI-1452 b、Kepler-138c 和 Kepler-138d 被发现其密度与其质量中有很大一部分由水构成的情形相一致 $^{[35][36]}$。此外，对高质量岩石行星 LHS 1140 b 的模型研究表明，其表面可能被一片深海所覆盖 $^{[37]}$。

尽管地球表面有 70.8\% 被水覆盖 $^{[38]}$，水在地球总质量中仅占约 0.05\%。在地外海洋世界中，海洋可能深且致密，以至于即便在高温条件下，巨大的压力也会使水转变为冰。在这类海洋的下部区域，数千巴的巨大压力可能促成奇异形态冰（如冰 V）地幔的形成 $^{[32]}$。这种冰并不一定像通常意义上的冰那样寒冷。若行星距离其恒星足够近，使水达到沸点，则水将进入超临界状态，不再具有清晰界定的表面 $^{[39]}$。即便是在较冷、以水为主的行星上，其大气层也可能比地球厚得多，并主要由水汽构成，从而产生极强的温室效应。此类行星必须足够小，无法保留厚重的氢和氦包层 $^{[40]}$，或者足够靠近其主星，以致这些轻元素被剥离 $^{[32]}$；否则，它们将演化成类似天王星和海王星那样的、更温暖版本的冰巨行星（此处尚需进一步证据）。
\subsection{历史}
20 世纪初的引力计算曾提出，木卫二（Europa）的组成可能富含水。1957 年，杰拉德·柯伊伯（Gerard Kuiper）通过地基观测证实了其表面主要由水冰构成。$^{[41]}$

在 20 世纪 70 年代行星探测任务之前，相关的重要理论基础工作已相继完成。1971 年，Lewis 指出，仅放射性衰变就可能足以在大型卫星中形成地下海洋，尤其是在存在氨（NH$_3$）的情况下。1979 年，Peale 和 Cassen 揭示了潮汐加热（亦称潮汐形变）在卫星演化与内部结构中的关键作用。$^{[3]}$
1992 年，人类首次确认探测到一颗系外行星。2003 年，Marc Kuchner 以及 2004 年 Alain Léger 等人提出，一小部分在雪线之外形成的冰质行星可能向内迁移至约 $1,\mathrm{AU}$ 的位置，其外层随后发生融化。$^{[42][43]}$

来自哈勃空间望远镜，以及先锋号、伽利略号、旅行者号、卡西尼–惠更斯号和新视野号等任务所积累的证据，强烈表明太阳系外侧的多个天体在隔热的冰壳之下蕴藏着内部液态水海洋。$^{[3][44]}$
与此同时，于 2009 年 3 月 7 日发射的开普勒空间望远镜已发现数千颗系外行星，其中约有 50 颗大小与地球相当，位于宜居带内或附近。$^{[45][46]}$

目前已探测到质量、尺寸和轨道各异的行星，这不仅展示了行星形成过程的多样性，也反映了行星在形成之后通过环恒星盘发生迁移的普遍现象。$^{[10]}$
截至 2025 年 10 月 30 日，已确认的系外行星共有 6,128 颗，分布在 4,584 个行星系统中，其中 1,017 个系统包含多于一颗行星。$^{[47]}$

2020 年 6 月，NASA 科学家基于数学建模研究报告称，银河系中拥有海洋的系外行星很可能相当普遍。$^{[11]}$

2022 年 8 月，距离地球约 99 光年的超级地球系外行星 TOI-1452 b 被凌日系外行星巡天卫星发现，该行星可能拥有深层海洋。$^{[35]}$
\subsection{形成}
\begin{figure}[ht]
\centering
\includegraphics[width=6cm]{./figures/00fa51d11bed7959.png}
\caption{阿塔卡马大型毫米/亚毫米阵列（ALMA）拍摄的 HL 金牛座（HL Tauri）原行星盘图像} \label{fig_Ocwo_5}
\end{figure}
在太阳系外侧形成的行星体，其初始状态通常类似彗星，按质量计大约由一半的水和一半的岩石组成，其整体密度低于岩质行星$^{[43]}$。在霜线附近形成的冰质行星和卫星应主要由 $\mathrm{H_2O}$ 和硅酸盐组成；而在更远区域形成的天体，则可以以水合物形式获得氨（$\mathrm{NH_3}$）和甲烷（$\mathrm{CH_4}$），并伴随一氧化碳（$\mathrm{CO}$）、氮气（$\mathrm{N_2}$）以及二氧化碳（$\mathrm{CO_2}$）$^{[48]}$。

在气态环星盘尚未消散之前形成的行星，会经历强烈的力矩作用，从而可能诱发其快速向内迁移进入宜居带，尤其对类地质量范围内的行星而言更为显著$^{[49][48]}$。由于水在岩浆中具有很高的溶解度，行星中相当大一部分水最初会被封存于地幔之中。随着行星逐渐冷却，地幔自下而上开始固结，大量水（约占地幔中水总量的 $60\%\sim99\%$）会发生析出，形成富含水汽的蒸汽大气层，并最终可能冷凝为海洋$^{[49]}$。海洋的形成需要行星内部发生分异过程，并且需要热i一个热源，这一热源可以来自放射性衰变、潮汐加热，或母体在早期阶段的高光度输出$^{[3]}$。然而，吸积完成之后的初始条件在理论上仍然缺乏充分约束。

在盘中外侧富水区域形成并随后向内迁移的行星，更有可能保有丰富的水储量$^{[50]}$。相反，在靠近其宿主恒星区域形成的行星，由于原始气体与尘埃盘的内侧通常被认为是高温且干燥的环境，因此其含水量往往较低。因而，如果在恒星附近发现一颗水世界，这将构成行星发生迁移并在原位之外形成的有力证据$^{[32]}$，因为在靠近恒星的区域缺乏足够的挥发性物质，难以支持原位形成$^{[2]}$。针对太阳系及系外行星系统形成过程的数值模拟表明，行星在形成过程中往往倾向于向内迁移（即朝向恒星）$^{[51][52][53]}$；在特定条件下，也可能发生向外迁移$^{[53]}$。向内迁移为冰质行星进入其冰层能够融化为液态水的轨道区域提供了可能性，从而使其转变为海洋行星。这一设想最早由 Marc Kuchner 于 2003 年在天文学文献中提出$^{[48]}$。
\subsection{结构}
冰质天体的内部结构通常通过对其整体密度、引力矩以及形状的测量来推断。确定天体的转动惯量有助于评估其是否经历过分异过程（即分离为岩石层与冰层）。在某些情况下，如果天体处于流体静力平衡状态（即在长时间尺度上表现为流体行为），则可以通过其形状或引力测量来反推转动惯量。尽管严格证明某一天体处于流体静力平衡状态极为困难，但通过综合利用形状与引力数据，仍可推导出其中的流体静力学贡献$^{[3]}$。用于探测内部海洋的具体技术包括磁感应测量、大地测量学、自转摆动、轴倾角、潮汐响应、雷达探测、成分证据以及表面形貌分析等方法$^{[3]}$。
\begin{figure}[ht]
\centering
\includegraphics[width=8cm]{./figures/1fed7f29913857c8.png}
\caption{木卫三（Ganymede）内部结构的艺术剖面示意图，其中液态水海洋被“夹”在两层冰层之间。各结构层按照真实比例绘制。} \label{fig_Ocwo_6}
\end{figure}
一种典型的冰质卫星由位于硅酸盐核心之上的水层构成。对于像土卫二（Enceladus）这样的小型卫星，海洋直接位于硅酸盐层之上、固态冰壳之下；而对于像木卫三（Ganymede）这样体量更大、富含冰的天体，其内部压力足够高，使得深处的冰转变为更高压相，从而有效形成一种“水夹层”结构，即液态海洋被夹在冰壳之间 $^{[3]}$。这两种情形之间的一个重要差异在于：在小型卫星中，海洋与硅酸盐直接接触，这可能为简单生命形式提供热液和化学能量及营养物质 $^{[3]}$。由于内部压力随深度变化，水世界模型中可能包含水的“蒸汽、液态、超流体、高压冰以及等离子体相”等多种状态 $^{[54]}$。其中部分固态水可能以冰 VII 的形式存在 $^{[55]}$。

维持地下海洋取决于内部加热速率与热量散失速率之间的对比，以及液体的凝固点 $^{[3]}$。因此，海洋的长期存在与潮汐加热机制密切相关。

较小的海洋行星通常具有更稀薄的大气层和更低的引力，因此液态水比在更大质量的海洋行星上更容易蒸发。数值模拟表明，质量小于一个地球质量的行星或卫星，仍可能在热液活动、放射性加热或潮汐形变的驱动下维持液态海洋 $^{[4]}$。当流体—岩石相互作用向深部脆性层传播较慢时，由蛇纹石化反应产生的热能可能成为小型海洋行星中热液活动的主要来源 $^{[4]}$。位于受潮汐形变影响的冰壳之下的全球性海洋，其动力学过程构成了一组重要而复杂的问题，目前才刚刚开始被系统研究。低温火山作用（cryovolcanism）发生的程度仍存在争议，因为水的密度比冰高约 8\%，在通常情况下不易喷发至表面 $^{[3]}$。尽管如此，旅行者 2 号、卡西尼–惠更斯号、伽利略号以及新视野号探测器的成像数据，已经在太阳系内多颗冰质天体表面揭示了低温火山活动的地貌特征。近期研究表明，在其内部海洋被表层冰覆盖的海洋行星上，低温火山作用可能与太阳系中的土卫二和木卫二（Europa）类似地发生 $^{[11][12]}$。

系外行星上的液态水海洋，其深度可能显著大于地球海洋的平均深度（约 3.7 km）$^{[56]}$。根据行星的引力和表面条件不同，系外行星的海洋深度可能达到地球的数十倍甚至数百倍。例如，对于一个表面温度约为 300 K 的行星，其液态水海洋深度可能在 30–500 km 之间，具体取决于行星的质量和组成 $^{[57]}$。
\subsection{大气模型}
